\chapter{Data Visualization and Experimentation (Iris)}

\section{Introduction}
In this session, we shift fully into an \textbf{Iris-only} workflow.
Recall from Session 8 that we already produced NumPy arrays (\texttt{X}, \texttt{y}) and class centroids.
Now, building on Session 8's outputs, we visualize classifier behavior and experiment results.

\begin{tcolorbox}[title=You will work in]
For this session, you will create or edit three files in the \texttt{session9/} folder:
\begin{itemize}
    \item \texttt{session9/iris\_eval.py} (helper functions for train/test split and accuracy)
    \item \texttt{session9/session9\_iris\_template.py} (main experiment runner)
    \item \texttt{session9/session9\_plot\_accuracy.py} (script for plotting)
\end{itemize}
\textbf{How to run:} From the project root: \\
\texttt{python3 session9/session9\_iris\_template.py} \\
\texttt{python3 session9/session9\_plot\_accuracy.py}
\end{tcolorbox}

You will work on two key tasks:
\begin{enumerate}\setlength{\itemsep}{2pt}
\item \textbf{Systematic parameter experiments:} run repeated experiments by changing a classifier setting (e.g., distance metric, threshold, or test size), and record the accuracy.
\item \textbf{Visualization:} plot how performance changes across settings, using box plots and line plots.
\end{enumerate}

\noindent
This tutorial is important because good data practice means:
\begin{itemize}\setlength{\itemsep}{0pt}
\item running \textbf{multiple trials} to reduce randomness and draw reliable conclusions,
\item saving results in a structured form (\textbf{list of dictionaries} $\rightarrow$ \textbf{DataFrame}),
\item visualizing distributions (box plot) and trends (line plot).
\end{itemize}

\subsection{Prerequisites}
You need:
\begin{itemize}\setlength{\itemsep}{0pt}
\item \texttt{numpy}, \texttt{pandas}, \texttt{matplotlib}, \texttt{openpyxl}
\item an \texttt{iris.csv} file
\end{itemize}

\noindent
Install if needed:
\begin{minted}[fontsize=\small, linenos, frame=lines]{bash}
pip install numpy pandas matplotlib openpyxl
\end{minted}

\newpage
\subsection{Calculating Accuracy for Iris Experiments}

In session 9, we combine prior ideas into an \textbf{experiment pipeline}:
\begin{itemize}\setlength{\itemsep}{0pt}
\item load Iris into NumPy arrays,
\item split the data into training and test sets,
\item run a classifier,
\item compute accuracy,
\item repeat under different settings,
\item store and visualize results.
\end{itemize}

\textbf{Important Parameters:}
\begin{itemize}\setlength{\itemsep}{0pt}
\item \texttt{test\_size}: The proportion of data used for evaluating the model. A larger test size leaves less data for training.
\item \texttt{seed}: Initializes the random number generator. Using a seed ensures our random data splits are reproducible.
\item \texttt{metric}: The distance metric used by the centroid classifier (\texttt{"l2"} is Euclidean distance, \texttt{"l1"} is Manhattan distance).
\end{itemize}

\subsubsection{Your Task: Reuse session 8 Modules}
Session 9 should import from Session 8 (e.g., \texttt{session8.iris\_session8\_utils} and \texttt{session8.iris\_session8\_models}) to preserve continuity.

\newpage
\subsection{Activity 1: Implement Train/Test Split and Accuracy}

\begin{tcolorbox}[title=Do: Create \texttt{session9/iris\_eval.py}]
We will write a custom split function using only NumPy:
\begin{minted}[fontsize=\small, linenos, breaklines=true, frame=lines]{python}
# session9/iris_eval.py
import numpy as np

def train_test_split(X, y, test_size=0.2, seed=1):
    """
    NumPy-only train/test split.
    """
    n = X.shape[0]
    rng = np.random.default_rng(seed)
    idx = np.arange(n)
    rng.shuffle(idx)

    test_n = int(n * test_size)
    test_idx = idx[:test_n]
    train_idx = idx[test_n:]

    return X[train_idx], X[test_idx], y[train_idx], y[test_idx]

def accuracy_score(y_true, y_pred):
    return float(np.mean(y_true == y_pred))
\end{minted}
\end{tcolorbox}

\newpage
\subsection{Activity 2: Refactoring for Flexibility}

A common beginner mistake is hardcoding values:
\begin{minted}[fontsize=\small, linenos, breaklines=true, frame=lines]{python}
# BAD PRACTICE (example)
test_size = 0.2
seed = 1
metric = "l2"
\end{minted}

This makes it hard to run systematic experiments. Instead, we'll create a function that takes these as arguments.

\begin{tcolorbox}[title=Do: Update \texttt{session9/session9\_iris\_template.py}]
Create the main experiment runner:
\begin{minted}[fontsize=\small, linenos, breaklines=true, frame=lines]{python}
# session9/session9_iris_template.py
import numpy as np
from session8.iris_session8_utils import load_iris_csv
from session9.iris_eval import train_test_split, accuracy_score
from session8.iris_session8_models import NumpyCentroidClassifier

def run_experiment(test_size=0.2, seed=1, metric="l2"):
    X, y = load_iris_csv("iris.csv", skip_header=1)
    X_train, X_test, y_train, y_test = train_test_split(X, y, test_size=test_size, seed=seed)

    model = NumpyCentroidClassifier(metric=metric)
    model.fit(X_train, y_train)
    y_pred = model.predict(X_test)

    acc = accuracy_score(y_test, y_pred)
    return acc

if __name__ == "__main__":
    acc = run_experiment(test_size=0.2, seed=1, metric="l2")
    print(f"One run | test_size=0.2 | seed=1 | metric=l2 | accuracy={acc:.3f}")
\end{minted}
\end{tcolorbox}

\begin{tcolorbox}[title=Checkpoint]
Run the file and ensure you get an accuracy output. Then, change the metric to \texttt{"l1"} and run it again to see if the accuracy changes.
\end{tcolorbox}

\newpage
\subsection{Running Multiple Experiments Dynamically}

Now we want to run multiple experiments with different parameter settings.

\begin{tcolorbox}[title=Do: Implement Looped Output]
Add this function inside \texttt{session9/session9\_iris\_template.py}:

\begin{minted}[fontsize=\small, linenos, breaklines=true, frame=lines]{python}
def run_single_setting_sweep(metrics, seeds, test_size=0.2):
    print(f"--- Sweep: test_size={test_size} ---")
    for metric in metrics:
        for seed in seeds:
            acc = run_experiment(test_size=test_size, seed=seed, metric=metric)
            print(f"metric={metric} | seed={seed} | acc={acc:.3f}")
\end{minted}

Then call it in the main block with \texttt{metrics=["l1", "l2"]} and \texttt{seeds=[1, 2, 3, 4, 5]}.
\end{tcolorbox}

\begin{tcolorbox}[title=Checkpoint]
Run a small sweep. You should see 10 lines of output representing the combination of 2 metrics and 5 seeds.
\end{tcolorbox}

\newpage
\subsection{Activity 3: Good Scientific Practice - Multiple Trials}

We run \textbf{multiple trials} for each configuration because randomness in the train/test split can cause accuracy to fluctuate. Averaging over many trials yields more reliable estimates.

\begin{tcolorbox}[title=Do: Implement \texttt{run\_multiple\_trials}]
\begin{minted}[fontsize=\small, linenos, breaklines=true, frame=lines]{python}
def run_multiple_trials(metrics, test_sizes, num_trials):
    """
    Returns a list of dicts:
    {'metric':..., 'test_size':..., 'trial':..., 'seed':..., 'accuracy':...}
    """
    results_list = []

    for metric in metrics:
        for test_size in test_sizes:
            print(f"\n[Config] metric={metric} | test_size={test_size}")
            for trial in range(1, num_trials + 1):
                seed = trial  # simple choice: seed=1..num_trials
                acc = run_experiment(test_size=test_size, seed=seed, metric=metric)

                print(f"  Trial {trial}/{num_trials} | seed={seed} | acc={acc:.3f}")

                run_data = {
                    "metric": metric,
                    "test_size": test_size,
                    "trial": trial,
                    "seed": seed,
                    "accuracy": acc
                }
                results_list.append(run_data)

    return results_list
\end{minted}
\end{tcolorbox}

\newpage
\subsection{Activity 4: Analyzing Results with Pandas}

DataFrames allow us to group, summarize, and easily save tabular data.

\begin{tcolorbox}[title=Do: Use Pandas to Process and Save Results]
At the top of \texttt{session9\_iris\_template.py}, add:
\begin{minted}[fontsize=\small, linenos, frame=lines]{python}
import pandas as pd
\end{minted}

In \texttt{main()}, after generating \texttt{all\_results\_data}:
\begin{minted}[fontsize=\small, linenos, breaklines=true, frame=lines]{python}
results_df = pd.DataFrame(all_results_data)
print(results_df.head())

output_filename = "iris_experiment_results.xlsx"
results_df.to_excel(output_filename, index=False)
print(f"Saved: {output_filename}")
\end{minted}
\end{tcolorbox}

\begin{tcolorbox}[title=Checkpoint]
Check your project directory for the newly created \texttt{iris\_experiment\_results.xlsx} file.
\end{tcolorbox}

\section{Troubleshooting Imports and Excel}
\begin{tcolorbox}[title=Common Mistakes \& Troubleshooting]
\begin{itemize}
    \item \textbf{ModuleNotFoundError:} Ensure you are running from the project root.
    \item \textbf{PermissionError (Excel):} If you get a "file in use" error when running your script, make sure you close the Excel file before trying to overwrite it!
\end{itemize}
\end{tcolorbox}

\newpage
\subsection{Activity 5: Analyzing and Visualizing Results}

\textbf{Box plots} show the distribution of accuracies, letting you see the median and variance. \textbf{Line plots} show the trend, such as how accuracy drops as test size increases.

\begin{tcolorbox}[title=Do: Visualization with Matplotlib]
Create \texttt{session9/session9\_plot\_accuracy.py}:
\begin{minted}[fontsize=\small, linenos, breaklines=true, frame=lines]{python}
import pandas as pd
import matplotlib.pyplot as plt

df = pd.read_excel("iris_experiment_results.xlsx")

# Box plot: accuracy distribution by metric
df.boxplot(column="accuracy", by="metric")
plt.title("Accuracy Distribution by Distance Metric")
plt.suptitle("")
plt.xlabel("Metric")
plt.ylabel("Accuracy")
plt.grid(True)
plt.savefig("session9_boxplot.png")
plt.show()

# Line plot: mean accuracy vs test_size
avg = df.groupby(["metric", "test_size"])["accuracy"].mean().reset_index()

plt.figure()
for metric in sorted(avg["metric"].unique()):
    sub = avg[avg["metric"] == metric]
    plt.plot(sub["test_size"], sub["accuracy"], marker="o", label=metric)

plt.title("Mean Accuracy vs Test Size")
plt.xlabel("Test Size")
plt.ylabel("Mean Accuracy")
plt.grid(True)
plt.legend()
plt.savefig("session9_lineplot.png")
plt.show()
\end{minted}
\end{tcolorbox}

\section{Summary and Reflection}
\textbf{Summary:}
\begin{itemize}
    \item We defined parameters like \texttt{test\_size} and \texttt{seed} to systematically test models.
    \item \textbf{Pandas} helped us collect dictionaries into DataFrames for easy aggregation and export.
    \item \textbf{Matplotlib} enabled us to visualize distributions and trends across metrics.
\end{itemize}

\textbf{Reflection Questions:}
\begin{enumerate}
    \item How does changing the \texttt{test\_size} impact the overall accuracy, according to your line plot?
    \item Why do we randomize seeds instead of using the same seed for every trial?
\end{enumerate}

\subsection{Submission}
Submit:
\begin{itemize}\setlength{\itemsep}{0pt}
\item \texttt{session9/session9\_iris\_template.py}
\item \texttt{session9/session9\_plot\_accuracy.py}
\item \texttt{session9/iris\_eval.py}
\item \texttt{iris\_experiment\_results.xlsx}
\item screenshots of plots (box plot + line plot)
\end{itemize}
